% Part: first-order-logic
% Chapter: semantics
% Section: intro-semantics

\documentclass[../../../include/open-logic-section]{subfiles}

\begin{document}

\olfileid{fol}{syn}{its}

\olsection{مقدمة}

إعطاء المعاني للتعبيرات من شأن الدلالة. والمفهوم المركزي في الدلالة
هو الإرضاء في !!a{structure}. تضفي !!^a{structure} معنى على اللبنات
الأساسية للغة: ف!!a{domain} مجموعة غير خالية من الأشياء. وتُؤوَّل
المكممات على أنها تتراوح على هذا المجال، وتُسنَد إلى ال!!{constant}s
عناصر من المجال، وتُسنَد إلى ال!!{function}s دوال من القوى المنتهية
المناسبة لل!!{domain} إلى المجال، وتُسنَد إلى ال!!{predicate}s علاقات على
ال!!{domain}. ويؤلف ال!!{domain}، مع الإسنادات إلى المعجم الأساسي،
!!a{structure}. وقد تظهر ال!!^{variable}s في ال!!{formula}s،
ولإعطائها دلالة يلزمنا أيضًا أن نسند إليها !!{element}s من
ال!!{domain}---وهذا إسناد للمتغيرات. وأخيرًا تجمع علاقة الإرضاء هذه
الأمور. فقد تُرضى !!^a{formula} في !!a{structure}~$\Struct{M}$ نسبة
إلى إسناد ل!!a{variable}~$s$، ويُكتب ذلك $\Sat{M}{!A}[s]$. وتُعرّف
هذه العلاقة أيضًا بالاستقراء على بنية~$!A$، باستخدام جداول الصدق
للروابط المنطقية لتعريف إرضاء $(!A \land !B)$، مثلًا، بدلالة إرضاء
$!A$ و~$!B$ (أو عدم إرضائهما). ثم يتبين أن إسناد ال!!{variable}s
غير ذي صلة إذا كانت ال!!{formula}~$!A$ !!a{sentence}، أي إذا لم تكن
فيها متغيرات حرة، وبذلك يمكننا الحديث ببساطة عن إرضاء ال!!{sentence}s
(أو عدم إرضائها) في ال!!{structure}s.

وعلى أساس علاقة الإرضاء $\Sat{M}{!A}$ لل!!{sentence}s يمكننا بعد
ذلك تعريف المفاهيم الدلالية الأساسية: الصلاحية والاستلزام وقابلية
الإرضاء. تكون !!^a{sentence} صالحة، $\Entails !A$، إذا أرضتها كل
!!{structure}. وتستلزمها مجموعة من ال!!{sentence}s،
$\Gamma \Entails !A$، إذا كانت كل !!{structure} ترضي جميع
ال!!{sentence}s في~$\Gamma$ ترضي أيضًا~$!A$. وتكون مجموعة من
ال!!{sentence}s قابلة للإرضاء إذا كانت !!{structure} ما ترضي جميع
ال!!{sentence}s فيها معًا. ولما كانت ال!!{formula}s معرّفة
استقرائيًا، وكان الإرضاء بدوره معرّفًا بالاستقراء على بنية
ال!!{formula}s، أمكننا استخدام الاستقراء لإثبات خصائص دلالتنا ولربط
المفاهيم الدلالية المعرّفة.

\end{document}
